\section{B3今週の標語書き初め公募会 講評}
\rightline{（文責:部活恐竜）}

\subsection{はじめに}
B3談話室では、「寝る前に二段ジャンプをしない」や「植木鉢だから一人になろう」といった今週の標語
が主に文責者によってほぼ毎週掲示されています。この文章は、そんな今週の標語を2026年1月1日
から5日にかけてB3全体から公募した際の講評となるものです。本講評が新入生の皆さんの標語観を
高める一助になれば幸いです。

\subsection{本文}
新春今週の標語書き初め公募会2026に応募していただいてありがとうございました。
告知もなく始めた今公募に34もの応募があったことを非常に嬉しく思っています。
こうした結果はひとえに、標語を求める現代人の深層心理の表れであると受け止めております。

当初は単純に最優秀賞作品だけ決めて標語に採用して終わろうと思っていたのですが、募集した以上
なんらかのレスポンスがあるのが筋だという意見を聞いて納得し、誠に勝手ながら応募していただいた
全ての標語に何らかのコメントをつけさせていただきたいと考えました。
応募していただいた標語に対する一方的な批判にはならないように細心の注意を払ったつもりなのでそ
の点は安心していただけますと幸いです。

\subsection{審査にあたって (読み飛ばしてよい!)}
審査結果を公表するにあたって、一定の審査基準を公表してそれに則ることにしました。
審査するのが私個人である以上、主観性や恣意性からはどうしたって逃げられませんが、他人に共有
する上では少なくともそういった感覚を納得できるような基準に変換するべきだと考えたからです。
標語という空間の機能を考えると、それは「それを見た人に強制的に読みのリソースを割かせること」に
あると思っています。それは一瞬チラ見する程度の出来事かもしれませんが、標語として提示されたそ
れが解読可能な文字である以上、それが有する意味を考える時間があるはずです。このとき、その文
字たちは自然言語上では不自然な接続だったとしても、断片だけでも意味が受け入れられて、言葉間
をつなげる意味を与えようとする試みが行われるでしょう。そうした処理の中で普段の言語使用では出
会うことの無いような新鮮な意味の側面を言葉が見せてくれることがあるはずです。
私にとっての良い標語とは、そのような可能性があるものです。

そのような良い標語を作るために必要な要素には以下のようなものがあると考えています。
審査基準はこれらの要素に準じます。

・言葉と言葉の距離
普段の言語使用で自然に成り立つような言葉の並びだと、標語として提示される面白みがありません。
かといって、全く関係ない二つの言葉を並べても、なんだこれと思われて終わってしまいます。その丁度
よい塩梅をついているかが言葉間の距離という観点です。

・私たちと言葉の距離
言葉間だけではなく、言葉ー私たちにも距離という観点は導入されます。つまり、普段見慣れているわ
けでもなく、全く知らない言葉でもない、久しぶりに聞いた・なんとなく想起できる言葉を使われると嬉しく
感じます。この基準に則ると季節性のある言葉をその季節に使っていた場合は言葉としての新鮮さは
少ないと判断されるでしょう。ハロウィンの時期にハロウィンの話をするのは日常会話でいいわけです。

・他者性
標語の論理構造が既知のものでないかという指標になります。距離で説明するなら、主張と私たちの距
離です。ついていける程度に耳慣れない論理、論法、意思は他人の話を聞く時のように新鮮に、魅力的
に映るはずです。

・字面
標語がどれくらい長いか、仮名と漢字のバランスがどのようになっているのかという観点です。長すぎる
標語は文字当たりの見てもらえる時間が少ないはず=読みのリソースが少ないはずですし、漢字が多
すぎると負荷が高くて同様の効果が生じるはずです。また、出てくる言葉の見た目が似ていると意味的
に違ってもニュアンスの距離が近づいて感じられてしまうこともあるでしょう。

以上のような基準に則って審査をさせていただきました。

\subsection{結果発表}

\subsubsection{最優秀賞}

・昼寝とモッツァレラ

単純な取り合わせですが今公募に応募していただいた全ての標語の中で二者の距離感が最も優れて
いると感じました。
全く一緒にいるのを見たことがない二つの語なのに、白、光、安寧、脱力、無防備、ご褒美といったイ
メージが重なり合って、まるでこれまでも使われてきた組み合わせかのように思えるくらいの納得感が
ありました。
そしてそれらは完全に重なり合っているわけでもなく、モッツァレラによって昼寝の窓からは暖かな光が
差し込み、昼寝によってモッツアレラは静かなキッチンのボウルの中に沈み、その両者を埋める景色や
文脈をゆっくりと、しかし着実に提示してくれます。
この二者の距離感を完全に信用してかそれ以上の物を付け足さない度胸にも胸を打たれました。
この標語を見て私たちが何をするべきかは語るまでもありません。
2026年は昼寝とモッツアレラから始めましょう。

\subsubsection{優秀賞}

・おにぎりをぶ厚くしよう

ぶ厚くするという一つの動詞でおにぎりに対して厚さというステータスを勝手に付与し、それを動かす旨
みをもたらしているのは非常に標語的な営みだと感じました。ただ二つの離れた要素をぶつけただけで
はなく、距離があるけれどもなんとなく温度感が似ているために標語全体のまとまりをよくしている語の
選択の丁寧さも高評価です。
使われている言葉はそれほど目新しいものではないけれども取り合わせによって新鮮さを創り出してい
るという点で標語の原義に立ち返っている作品だと言えると思います。

・パリと結ばれたい!

パリ、結ばれるという言葉の選択は距離感が近そうでヒヤッとさせられますが、それを助詞の「と」によっ
て想定されていた意味を外れて既視感を回避し、新たなニュアンスを発生させていて標語的に評価でき
る手法だなと感じました。
文末の願望や!によって生じる勢いも、こうした「と」によってパリー結ばれる間に無理矢理生じた意味
の経由地点の強引さを加速させていて、認知ののどごしが気持ちいい標語だといえます。

・つるつるの壺

文章でも取り合わせでもない挑戦作ですが、純粋にいいなと思いました。
最優秀賞が昼寝とモッツアレラの関係を観せてくれる標語だとするならば、この標語はこの標語を見て
いる我々とつるつるの壺の距離を観せてくれている作品でしょう。
つるつるの壺という言葉と対峙したとき、私たちは何を思うでしょうか。まず、語感が良い。これが素晴ら
しい。語感が良いから何度も頭の中で唱えているうちに頭のなかで壺が造形される。でも思ったよりも
具体的な像が結ばれない。しかし、ちょっと考えてたら出来そうだから造形し続ける。つるつるのという
オーダーに従って表面を整備する。そして、頭の中でひとしきりそれっぽいものを作った後に結局なん
だったんだって思う。最終的にはそう思うのですが、その過程には眼前に立ち現れた立体的なつるつる
の壺が存在するはずなのです。
このような我々のつるつるの壺に対する歩みそのものが標語的だと言えるのではないでしょうか。
そうした効果が作者の意図的なものなのかはさておき、作者が選択した壺という対象が想像をする上で
も、その想像をなんだったんだと棄却する上でも魅力的な対象だったことは確かです。
類似標語に「スポーツの秋」(2024 3/18~3/24)がありますが、この標語はよりゆっくりと我々と提示され
た言葉の距離を埋めてくれています。標語の地平もここまで拓けたのかと驚きました。

\subsubsection{佳作}

・チョロギの持つエネルギーを最大化させよう
チョロギに勝手にエネルギーという概念を導入してそれをどうこうしようとしているのは、おに厚に類する
良さがあったのですが、肝心のチョロギが普段の自分からの距離が遠すぎるせいであれって何だっけと
いう考えがワンクッション挟まってしまったのが惜しかったです。これは私のチョロギに対する感度が低
かったことにも責任があります。精進。

・鏡餅を回転体としてとらえよう
普段の言語使用からは距離がある上で素朴さがある回転体という言葉の選択は評価できますが、鏡餅
が回転体だという主張が納得できすぎてしまいました。納得に使うカロリーが少なかったせいで読みの
リソースが余ってしまい、少し物足りなさを感じてしまったのが惜しく感じました。

・擬宝珠をメガ小籠包に取り替えよう
語の新規性や動詞のちょうどよさはかなり評価できますが、造語が新しすぎてイメージの起点が作りず
らいのが難しかったです。これは私の最近の標語でも課題とされている部分なので傾向をかなりおさえ
て下さっている方の作品かもしれません。そうであれば課題を抱えた標語を学習させてしまって申し訳
ないです。

・暴言をジェルにして飲み込む
飲み込むを起点にして暴言→飲み込むとジェル→飲み込むという二つの文脈を連れてきて重ねようとし
ているやり方は非常に標語的なだなと感じました。しかし、ジェル→飲み込むの対象として想起されるイ
メージと暴言のイメージが近すぎてジェルと暴言の距離が近づいてしまい、暴言をジェルにする面白み
が減ってしまっているのがそれだけに惜しく感じられました。

・琵琶湖の底でツチノコを養殖しよう
琵琶湖の底という存在は確かだけど考えたことのない場所をついてくれたのは嬉しかったです。一方
で、ツチノコの養殖はツチノコにするべきではないことランキングの上位に存在してしまうので距離が近
いように感じられました。

・もろみとトリカブトをデジクロスさせてPPAPに取り込もう
語の距離感や字面のバランスは素晴らしいのですが、デジクロスという言葉が指し示すものが取り合わ
せの行為そのものすぎてこれらの語が集まったことに不必要な納得性が生じてしまっているのがとても
惜しかったです。

・ら ゙を使う

謎の目的意識を有している点は自分とは遠い他者性があって素晴らしいのですが、「ら ̈」という普段使
わない語と使うという動作の距離が完全に反対すぎて二者の関係にすぐに思い至れてしまうのが惜しく
感じられました。

・雑煮を蛇腹折にしよう
今までに述べた審査基準は概ねクリアしてそうなのですが、イメージするカロリーが高すぎて心にささり
きりませんでした。二者の距離が遠すぎるということなのかもしれませんし、私の標語力不足かもしれま
せん。精進します。

・来年の目標を考える
標語じゃなくて目標という語を選択しているのは標語を見ている我々との距離が遠くて良かったのです
が、来年というアイテムと新年と言う時期にいる我々との距離がかなり近く、それだけに普遍的なもので
あるべきである標語が非常に個人的なものに収まってしまったように感じられました。

・心を亀裂骨折までで留める
亀裂骨折という語の選択は我々との距離がいい具合に遠くて素晴らしいのですが、心と骨折の距離が
近すぎてしまってまともなテーマを感じられてしまうのが標語的には相性がよくないように感じられまし
た。

・乾物に浮気しよう
あえて何から浮気するのかについての言及を避けることで文章量以上の広がりがあったのは非常に標
語的な理念に沿っていて良かったです。ただ乾物と浮気の距離が思いのほか近く、文面上に提示され
ている一本の面白がり方以外を考えにくかったのが惜しく感じました。

・凧で多幸感を得よう
二者の意味的な距離はいい塩梅なのですが、音の距離が近すぎてただそれだけに見えてしまいかね
ないのが惜しく感じました。

・守り抜こう、この対称性を
見たことのある文章や文脈に異物を代入するやり方はこれまでの標語でも行われていることで、そうし
たやり方を自分なりに見出して踏襲している点はこの公募会というコンセプトに適っていて素晴らしいで
す。ただし、代入された対称性という言葉と元の文脈で入り得る言葉との距離が近すぎて異物感が薄れ
てしまっているのが惜しく感じられました。

・最近は歳で、ピクミンとミニドラの区別もつかなくなってきてしまいました。
こちらの標語でも代入が行われている上で元の文脈との解離がなされていて素晴らしいのですが、字
面的にもピクミンとミニドラが似たようなイメージで捉えられてしまい、区別がつかないことに対してある
種の当然さを感じてしまう隙があったのが惜しかったです。

・数学的帰納法で未来の標語を予測しよう
文章全体の「したとて」感は素晴らしいのですが、数学的帰納法自体が未来や予測との関連が強すぎ
て文面以上の広がりを感じるのが難しかったです。標語という単語が数学的帰納法についての文章に
代入されたと考えてもいいのかもしれませんが、その場合は代入する単語が自己言及的過ぎて悪目立
ちしてしまうように感じました。

・草餅を雑煮に入れて雑草煮にせよ
イメージしやすい造語や謎の責務感などは評価できるポイントですが、雑草煮ができる論理が分かりや
すすぎてしまって読みのリソースが余ってしまったように感じられました。

・眼鏡に蛍光ステッカーを貼ってしまおう
文末によって貼るという行為を阻害する謎の力が見えてくるのは標語的に評価できるのですが、眼鏡と
蛍光ステッカーという選択が二者の取り合わせとしても文脈としても旨みを足しきれていないように感じ
られてしまいました。

・天井の穴シェアリングサービス開始!
全く意義の無い行動を堂々としているのは不気味な他者という感じで良かったのですが、意義の無さに
たどり着く論理が見えやすすぎて勢いが空回りしているようにも感じられてしまいました。

・ヒョウモンダコとも仲良くしよう
助詞の「も」によってヒョウモンダコ以外の仲良い存在やヒョウモンダコそのものに対する見下しといった
言外の情報量を高めている点が評価できます。しかし、こういった評価点は「も」のちからによるもので
あり、独自性が現れるヒョウモンダコや仲良くするといった言葉の選択に旨みが感じにくかったのが惜し
かったです。

・人間とは、回らないコマである。
人間とは、で広げた大風呂敷の中心に脱力した言葉が存在するのは好みですが、少しかっこよすぎる
ように思えました。かっこよすぎるせいで頭のどこかが真剣になってしまって冷めた自分が心に現われ
てしまったのが惜しかったです。

・人間に本当の食物連鎖を理解らせる
言い切りによっても増している目的意識の強さは他者性があって良いのですが、人類に対する強い憎
しみが原液で垣間見えてしまっているのが惜しく感じられました。標語に感情は不要。

・エクゾディアでババ抜きしよう
イメージのトンチキさは評価できるのですが、エグゾディアという固有名詞が出て来たときにそこから離
れることを期待してしまう中で、着地点のババ抜きという行為がエグゾディアでやってはいけないことに
見えてしまって関連を感じすぎてしまいました。

・爆破
文章ですらないという点では優秀賞である「つる壺」と同じなのですが、爆破というものがオチすぎてい
てそこから先に何かを考えることが難しくこの評価となりました。

・破戒
文章ですらないという点では優秀賞である「つる壺」と同じなのですが、破戒というものがオチすぎてい
てそこから先に何かを考えることが難しくこの評価となりました。

・パキすぎてポキ!
下の欄に送ってやろうか悩みましたが、ポキがパキよりも上だという謎の上下関係が見えて嬉しくなっ
たので論内とします。

\subsubsection{論外}

・バカたれ
単純な暴言なのかバカであれという内容なのかどちらか分かりませんがどちらにせよ論外です。

・バカがよ
単純な暴言なのか酷い国家の替え歌なのかどちらか分かりませんがどちらにせよ論外です。

・泣きアニメ
奥座敷先輩、好きな言葉をとりあえず送ってくるのはやめてください。

・トロリーバス 猛暑で溶けて トロトロリーバス
かなりくだらなくて素晴らしいです。単純なインパクトだけではなく、猛暑で溶ける妥当性がまったく存在
しない点も好印象です。惜しむらくは、くだらなすぎることです。くだらなすぎて、こんなものを掲示してい
ては標語に対する信頼性を失ってしまいます。泣く泣くですが破門です。

\subsection{おわりに}
ここまで読んでいただいて本当にありがとうございました。
私以外の人間が標語というものを解釈して再構成したものがこれだけの数見れたことに感動していま
す。中でも優秀賞以上の作品には、こうした標語が私以外の人間から作られたことに悔しさすら覚えま
す。こうした営みの中で皆さんが許容できる言葉の距離が長くなって、魅力的な標語がどんどん作られ
ていけば、そして日常会話に支障が生じていけばそれ以上の喜びはありません。
また、審査基準を作成するにあたって、自分が標語についてどのように考えているのか改めて考える機
会を作っていただけたことに感謝します。このように審査を公開した以上、私の今までの標語が同様の
基準で眼差される機会も生まれると思います。そのときはどうか自分なりに厳しい言葉を投げかけてく
ださい。自分なりの標語の分派をつくっていただいてもかまいません。今回はこのように私一人が審査
を行いましたが、このように標語というものが私一人の思想によって規定されてしまうことを本意には考
えていません。私たちで共に新たな標語の地平を切り拓いていきましょう。